\section{Introduction}
\label{sec:introduction}

Modern physics describes empirically successful regularities using several effective
frameworks (relativity, quantum theory, gauge fields), but their mutual conceptual
compatibility remains debated, and ``quantization'' is typically imposed as a postulate.
This paper explores an alternative \emph{substrate-first} strategy: assume an ontically real,
deterministic wave medium, and treat familiar field theories as coarse-grained descriptions
of constrained wave dynamics.

\paragraph{Substrate hypothesis.}
We assume a three-dimensional elastic medium (``brane'') embedded in $\R^4$, described by
an embedding map
\begin{equation}
  \vecX:\Omega\times\R \to \R^4, \qquad x\mapsto \vecX(x,t), \qquad \Omega\subset\R^3.
  \label{eq:embedding-map}
\end{equation}
Time $t$ is an external evolution parameter (not a geometric coordinate).
The brane is assumed isotropic in its reference coordinates $x$.

\paragraph{Motivation for geometric phase as a gauge diagnostic.}
Multi-component waves generically exhibit geometric phases under adiabatic or cyclic transport.
In optics this is known as the Pancharatnam--Berry phase \citep{Pancharatnam1956,Berry1984,Cisowski2022},
and non-Abelian generalizations arise for degenerate subspaces \citep{WilczekZee1984}.
The central technical idea in this paper is to treat the geometric phase of a \emph{narrowband, polarized}
substrate wave as a \emph{diagnostic} for candidate gauge structures, rather than postulating
a $U(1)$ field a priori.

\paragraph{Relation to analogue/emergent gravity programs.}
Lorentzian kinematics and effective metrics can emerge in wave systems via linearization and
hydrodynamic limits \citep{Barcelo2011AnalogueGravity}. We use this perspective to motivate
an \emph{effective} light cone, while keeping the microscopic dynamics strictly Newtonian in $t$.

\paragraph{Related work.}
The present approach intersects three bodies of work:
(i) geometric phase and gauge structures in wave systems \citep{Berry1984,Pancharatnam1956,WilczekZee1984,Cisowski2022,Bliokh2015SpinOrbit},
(ii) continuum mechanics and isotropic hyperelasticity as rotationally invariant substrates \citep{Ogden1984,Holzapfel2000,DoCarmo1992},
and (iii) emergent/analogue gravity ideas \citep{Barcelo2011AnalogueGravity,Sakharov1967VacuumGravity,Volovik2003HeliumDroplet,Jacobson2004EinsteinAether}.
Brane-world models in high-energy theory use ``branes'' in a different (field-theoretic) sense \citep{RandallSundrum1999};
here the term is used literally as an elastic medium.

\paragraph{Paper organization.}
\Cref{sec:scope} states scope, contributions, and non-claims.
\Cref{sec:continuum-model} formulates the continuum substrate model and stresses its rotational invariance.
\Cref{sec:emergent-relativity} discusses the linear-wave regime and an emergent effective light cone.
\Cref{sec:geophase} develops geometric-phase diagnostics and their gauge freedom.
\Cref{sec:localized} outlines localized modes and how to pose them as self-adjoint eigenproblems.
\Cref{sec:validation} summarizes validation targets and falsifiable checks.
Implementation details are in \Cref{app:discrete}.
